\section{Cycle--Overlap Visibility Theorem}

Let \(G\) be a finite graph with maximum degree \(\Delta\).
Fix \(R \in \mathbb{N}\).

\subsection{Cycle--Overlap Rank}

Let \(\mathcal{C}_R(G)\) denote the set of cycles contained in radius-\(R\) balls.
Define the cycle-overlap rank

\[
COR_R(G) =
\dim_{\mathbb{F}_2}
\mathrm{span}\{\chi_C : C \in \mathcal{C}_R(G)\}
\subseteq \mathbb{F}_2^{E(G)} .
\]

\subsection{Transport Signature}

For a vertex \(v\),

\[
\Theta_R(v) =
\big(
\deg(v),
\{\tau_R(v,u)\}_{u \in N(v)}
\big)
\]

where

\[
\tau_R(v,u) =
\#\{C \in \mathcal{C}_R(G) : (v,u) \in C\}.
\]

\subsection{Rigidity Lemma}

There exists a constant \(T(\Delta,R)\) such that

\[
COR_R(G) \ge T(\Delta,R)
\]

implies

\[
\exists u,v :
\Theta_R(u) \neq \Theta_R(v).
\]

\subsection{WL\(^2\) Consequence}

Vertices with different transport signatures receive different
WL\(^2\) colors.

\[
\Theta_R(u) \neq \Theta_R(v)
\Rightarrow
\chi_{WL^2}(u) \neq \chi_{WL^2}(v).
\]

\subsection{FO\(^3\) Consequence}

WL\(^2\) distinguishability implies FO\(^3\) type difference via
the \(3\)-pebble Ehrenfeucht--Fra\"iss\'e game.

\[
\chi_{WL^2}(u) \neq \chi_{WL^2}(v)
\Rightarrow
tp_3(u) \neq tp_3(v).
\]

\subsection{Cycle--Overlap Visibility}

Combining the steps yields

\[
COR_R(G)
\Rightarrow
\Theta_R \text{ diversity}
\Rightarrow
WL^2
\Rightarrow
FO^3.
\]

\begin{theorem}[Cycle--Overlap Visibility]
For bounded-degree graphs there exists \(T(\Delta,R)\) such that

\[
COR_R(G) \ge T(\Delta,R)
\]

implies the existence of vertices with distinct FO\(^3\) types.
\end{theorem}
